\section{Há proteção intelectual do(s) produto(s) ou processo(s) mais famosos da empresa? Qual o tipo?}

A \textbf{NVIDIA Corporation} protege suas inovações e produtos através de diversos mecanismos de propriedade intelectual, conforme descrito em seu código de conduta:

\begin{itemize}
    \item \textbf{Patentes}: A NVIDIA possui patentes que protegem suas principais tecnologias, incluindo arquiteturas de GPUs e DPUs. Essas patentes garantem exclusividade e limitam a concorrência, oferecendo à empresa direitos exclusivos sobre o uso e licenciamento de suas inovações.

    \item \textbf{Direitos Autorais (Copyright)}: Os softwares e algoritmos desenvolvidos pela NVIDIA, como drivers e bibliotecas de IA, são protegidos por direitos autorais, impedindo cópia ou distribuição não autorizada.

    \item \textbf{Marcas Registradas (Trademarks)}: A NVIDIA mantém várias marcas registradas, como “NVIDIA” e “GeForce,” que protegem a identidade dos produtos e impedem que terceiros usem essas marcas sem autorização.

    \item \textbf{Segredos Industriais (Trade Secrets)}: Para proteger informações confidenciais, como métodos de fabricação e algoritmos proprietários, a NVIDIA utiliza segredos industriais, reforçados por acordos de confidencialidade com colaboradores e parceiros \cite{nvidia_code_conduct}.
\end{itemize}